% !TeX program = pdflatex
\documentclass[12pt,a4paper]{report}   % 12pt вместо 14pt

\usepackage[T2A]{fontenc}
\usepackage[utf8]{inputenc}
\usepackage[russian]{babel}
\usepackage{geometry}
\usepackage{epstopdf}
\usepackage{setspace}
\usepackage{graphicx}
\usepackage{amsmath,amssymb,amsthm,mathtools}
\usepackage{booktabs}
\usepackage{float}
\usepackage[unicode]{hyperref}
\usepackage{indentfirst}
\usepackage{enumitem}
\usepackage{titlesec}
\usepackage{caption}
\usepackage{longtable}
\usepackage{xcolor}
\usepackage{tempora}
\usepackage{tocloft}

% Поля: левое 30 мм, правое 15 мм, верх/низ 20 мм
\geometry{left=30mm,right=15mm,top=20mm,bottom=20mm}
\onehalfspacing                     % полуторный интервал
\setlength{\parindent}{1.25cm}

\hypersetup{
    colorlinks=true,
    linkcolor=black,
    urlcolor=blue,
    citecolor=black
}

% Настройка заголовков глав и разделов
\titleformat{\chapter}[block]
    {\normalfont\bfseries\centering}
    {\thechapter.}
    {1em}
    {}
\titlespacing*{\chapter}{0pt}{-35pt}{12pt}

\titleformat{\section}
    {\normalfont\bfseries}
    {\thesection.}
    {1em}
    {}

\titleformat{\subsection}
    {\normalfont\bfseries}
    {\thesubsection.}
    {1em}
    {}

\titleformat{name=\subsection,numberless}
    {\normalfont\bfseries}
    {}
    {0pt}
    {}

\captionsetup{labelsep=period}     % точка после номера рисунка/таблицы

% Оглавление
\renewcommand{\contentsname}{Содержание}
\setlength{\cftbeforetoctitleskip}{-2.5em}
\setlength{\cftaftertoctitleskip}{1.5em}
\renewcommand{\cfttoctitlefont}{\bfseries\large}
\renewcommand{\cftaftertoctitle}{}

% Список литературы
\addto\captionsrussian{\renewcommand{\bibname}{Список литературы}}

\begin{document}
    \begin{titlepage}
        \begin{center}
        \vspace*{0.2cm}
        
\includegraphics[width=6.6cm]{images/msu_logo.png}
        \vspace{0.2cm}

        {\fontsize{12pt}{14pt}\selectfont
        МОСКОВСКИЙ ГОСУДАРСТВЕННЫЙ УНИВЕРСИТЕТ ИМЕНИ М.\,В.~ЛОМОНОСОВА\par}
        {\fontsize{12pt}{14pt}\selectfont
        ФАКУЛЬТЕТ ВЫЧИСЛИТЕЛЬНОЙ МАТЕМАТИКИ И КИБЕРНЕТИКИ\par}
        {\fontsize{12pt}{14pt}\selectfont
        КАФЕДРА ВЫЧИСЛИТЕЛЬНЫХ МЕТОДОВ\par}

        \vspace{2.5cm}

        {\fontsize{15pt}{17pt}\selectfont
        Козлов Кирилл Андреевич\par}

        \vspace{0.25cm}

        {\bfseries\fontsize{16pt}{18pt}\selectfont
        Численное решение модели реального делового цикла\par
        для закрытой экономики\par}

        \vspace{1.6cm}

        {\fontsize{16pt}{18pt}\selectfont
        Численное решение модели реального делового цикла для закрытой экономики\par}
        \end{center}

        \vspace{1.6cm}
        \noindent
        \hspace*{6.0cm}
        \begin{minipage}{0.43\textwidth}
        {\bfseries\fontsize{15pt}{17pt}\selectfont
        Научный руководитель:\par}
        \vspace{0.15cm}
        {\fontsize{14pt}{16pt}\selectfont
        профессор, д.ф.-м.н.\par
        Богомолов Сергей Владимирович\par}
        \end{minipage}
        \vfill
        \begin{center}
        {\fontsize{15pt}{17pt}\selectfont Москва, 2026}
        \end{center}
        \vspace*{-0.8cm}
    \end{titlepage}

    % Начинаем нумерацию страниц с оглавления (стр. 2)
    \setcounter{page}{2}
    \pagenumbering{arabic}

    \tableofcontents
    \clearpage


    % --- Введение ---
    \chapter*{Введение}
        \addcontentsline{toc}{chapter}{Введение}

        Модели динамического общего равновесия занимают важное место в современной макроэкономической теории. Они позволяют исследовать совместную динамику выпуска, потребления, капитала и труда, а также анализировать механизм межвременного выбора экономических агентов. Одной из наиболее известных конструкций такого типа является модель реального делового цикла (Real Business Cycle, RBC), в которой колебания основных макроэкономических переменных объясняются реакцией экономики на реальные шоки при рациональном поведении домохозяйств и фирм.

        Актуальность темы обусловлена тем, что даже в детерминированной постановке модель реального делового цикла приводит к нелинейной системе разностных уравнений, для которой аналитическое построение траектории, вообще говоря, невозможно. В связи с этим ключевую роль приобретают численные методы. Особый интерес представляет случай, когда в уравнении Эйлера доходность капитала зависит от величин следующего периода. Тогда переход по потреблению определяется неявно, и на каждом шаге численного алгоритма требуется решать дополнительное нелинейное уравнение.

        Целью выпускной квалификационной работы является построение и численное исследование детерминированной модели реального делового цикла для закрытой экономики, нахождение ее стационарного состояния и реализация численного решения методом стрельбы. А также численный анализ модели с учётом шоков, которые могут произойти в экономике. 

        Для достижения поставленной цели необходимо решить следующие задачи:
        \begin{enumerate}[label=\arabic*)]
            \item выписать систему уравнений модели и свести ее к замкнутой динамической форме;
            \item аналитически получить формулы для стационарного состояния;
            \item рассмотреть существующие численные методы решения задач оптимального управления в макроэкономике;
            \item реализовать численный алгоритм решения методом стрельбы;
            \item построить переходные траектории основных макроэкономических переменных;
            \item исследовать чувствительность решения к выбору параметров модели.
        \end{enumerate}

        % Объектом исследования является детерминированная макроэкономическая модель закрытой экономики типа RBC. Предметом исследования выступают стационарное состояние модели, переходная динамика капитала, потребления, выпуска и труда, а также численные методы построения устойчивой траектории.

        % Работа состоит из введения, пяти глав, заключения и списка литературы. В первой главе формально ставится задача и выписывается система уравнений модели. Во второй главе рассматриваются существующие численные подходы к решению задач оптимального управления в экономике. В третьей главе проводится аналитическое исследование стационарного состояния и выводится двумерная динамическая система. В четвертой главе описывается программная реализация численного метода и приводятся результаты вычислительных экспериментов. В пятой главе формулируются основные выводы исследования.

    % --- Глава 1 ---
    \chapter{Постановка задачи}
        \section{Описание детерминированной модели}
            Рассматривается детерминированная модель реального делового цикла в дискретном времени. Динамика экономики описывается следующей системой уравнений:
            \begin{equation}
            \begin{cases}
            \dfrac{C_{t+1}}{C_t} = \beta\bigl(1+R_{t+1} - \delta\bigr) \\[6pt]
            K_{t+1} = I_t + (1 - \delta)K_t \\[4pt]
            Y_t = A K_t^{\alpha} L_t^{1-\alpha} \\[4pt]
            L_t = \left(\dfrac{w_t}{\phi C_t}\right)^{\frac{1}{\psi}} \\[6pt]
            R_t = \alpha \dfrac{Y_t}{K_t}\\[6pt]
            w_t = (1-\alpha)\dfrac{Y_t}{L_t} \\[6pt]
            Y_t = C_t + I_t
            \end{cases}
            \end{equation}

            Здесь $C_t$ обозначает потребление в момент времени $t$, $K_t$ --- объём капитала, $L_t$ --- предложение труда, $Y_t$ --- выпуск(производство), $I_t$ --- инвестиции, $R_t$ --- доходность капитала, $w_t$ --- заработную плату. Параметр $\beta\in(0,1)$ является коэффициентом межвременного дисконтирования, $\alpha\in(0,1)$ --- эластичностью выпуска по капиталу, $A>0$ --- уровнем совокупной факторной производительности, $\delta\in(0,1)$ --- нормой выбытия капитала(ставкой амортизации), $\phi>0$ --- параметром дизполезности труда(фактор ненависти к труду), а $\psi>0$ --- параметром, связанным с эластичностью предложения труда по Фришу.

            В рассматриваемой модели реального делового цикла все переменные являются функциями дискретного времени $t = 0,1,2,\dots$ Ключевым \textbf{переменным состояния} выступает запас капитала $K_t$, динамика которого определяется законом накопления и полностью характеризует положение экономики в каждый момент времени. \textbf{Переменными управления} служат потребление $C_t$ и предложение труда $L_t$, которые выбираются домохозяйством для максимизации своей межвременной полезности. Остальные макроэкономические показатели — выпуск $Y_t$, доходность капитала $R_t$, заработная плата $w_t$ и инвестиции $I_t$ — являются \textbf{эндогенными} и однозначно выражаются через переменные состояния и управления с помощью технологических ограничений и условий оптимальности. В частности, выпуск задаётся производственной функцией Кобба–Дугласа $Y_t = A K_t^{\alpha} L_t^{1-\alpha}$, доходность капитала и ставка заработной платы определяются предельными продуктами соответствующих факторов: $R_t = \alpha Y_t/K_t$, $w_t = (1-\alpha)Y_t/L_t$, а инвестиции вычисляются из тождества $I_t = Y_t - C_t$. Таким образом, вся динамика модели сводится к эволюции пары $(K_t, C_t)$ при эндогенном определении труда $L_t$ из условия внутрипериодного выбора.

        \section{Формальная постановка задачи}

            Требуется исследовать динамику модели при заданных параметрах $\beta$, $\alpha$, $A$, $\delta$, $\phi$, $\psi$ и начальном значении капитала $K_0$. Для этого необходимо:

            \begin{enumerate}[label=\arabic*)]
                \item выразить все эндогенные переменные, включая труд $L_t$, выпуск $Y_t$, доходность капитала $R_t$ и заработную плату $w_t$, через текущие значения единственной переменной состояния $K_t$ и переменной управления $C_t$ (потребление);
                \item получить стационарное состояние $(K^*, C^*, L^*, Y^*, R^*, w^*)$ модели;
                \item построить переходную траекторию, сходящуюся к стационарной точке, путём подбора начального значения потребления $C_0$ при фиксированном $K_0$;
                \item реализовать численный алгоритм, позволяющий находить такое $C_0$, чтобы траектория на заданном горизонте выходила на стационарное состояние;
                \item провести вычислительные эксперименты для различных начальных условий $K_0$ (ниже и выше стационарного уровня).
            \end{enumerate}

        \section{Сведение системы к динамической форме}

            Из уравнений для заработной платы, предложения труда и производственной функции Кобба--Дугласа можно получить явное выражение для $L_t$ через $K_t$ и $C_t$. Исходные соотношения имеют вид:
            \begin{align}
                w_t &= (1-\alpha)\frac{Y_t}{L_t}, \label{eq:wage} \\
                Y_t &= A K_t^{\alpha} L_t^{1-\alpha}, \label{eq:cobb-douglas} \\
                L_t &= \left(\frac{w_t}{\phi C_t}\right)^{\frac{1}{\psi}}. \label{eq:labor-supply}
            \end{align}

            Подставляя (\ref{eq:wage}) в (\ref{eq:labor-supply}), получаем:
            \[
            L_t = \left( \frac{(1-\alpha)Y_t/(L_t)}{\phi C_t} \right)^{\frac{1}{\psi}}
                = \left( \frac{(1-\alpha)Y_t}{\phi C_t L_t} \right)^{\frac{1}{\psi}}.
            \]

            Заменим выпуск $Y_t$ согласно производственной функции (\ref{eq:cobb-douglas}):
            \[
            L_t = \left( \frac{(1-\alpha)A K_t^{\alpha} L_t^{1-\alpha}}{\phi C_t L_t} \right)^{\frac{1}{\psi}}
                = \left( \frac{(1-\alpha)A K_t^{\alpha} L_t^{-\alpha}}{\phi C_t} \right)^{\frac{1}{\psi}}.
            \]

            Возведём обе части уравнения в степень $\psi$:
            \[
            L_t^{\psi} = \frac{(1-\alpha)A K_t^{\alpha} L_t^{-\alpha}}{\phi C_t}.
            \]

            Умножая левую и правую части на $L_t^{\alpha}$, окончательно находим:
            \[
            L_t^{\psi+\alpha} = \frac{(1-\alpha)A}{\phi} \cdot \frac{K_t^{\alpha}}{C_t}. \label{eq: L-explicit}
            \]

            Таким образом, величина трудовых затрат $L_t$ однозначно выражается через текущие значения капитала $K_t$ и потребления $C_t$:
            \[
            L_t = \left( \frac{(1-\alpha)A}{\phi} \cdot \frac{K_t^{\alpha}}{C_t} \right)^{\frac{1}{\psi+\alpha}}.
            \]


            Поскольку труд $L_t$ уже выражен через $K_t$ и $C_t$, выпуск $Y_t$ и доходность капитала $R_t$ также становятся функциями только этих двух переменных. Подставляя найденное выражение для $L_t$ в производственную функцию Кобба--Дугласа, получаем
            \begin{equation}
            Y_t = A K_t^{\alpha} L_t^{1-\alpha} = A K_t^{\alpha} 
            \left( \frac{(1-\alpha)A}{\phi} \cdot \frac{K_t^{\alpha}}{C_t} \right)^{\frac{1-\alpha}{\psi+\alpha}}.
            \end{equation}
            Однако для практических вычислений удобнее сохранить промежуточное вычисление $L_t$ и затем использовать стандартные формулы:
            \begin{equation}
            Y_t = A K_t^{\alpha} L_t^{1-\alpha}, \qquad
            R_t = \alpha \frac{Y_t}{K_t}.
            \end{equation}
            Таким образом, зная $K_t$ и $C_t$, можно последовательно определить $L_t$, затем $Y_t$ и $R_t$, после чего перейти к следующему периоду.

            После подстановки выражений для выпуска $Y_t$ и инвестиций $I_t = Y_t - C_t$ закон движения капитала приобретает явную форму:
            \begin{equation}
            K_{t+1} = Y_t - C_t + (1-\delta)K_t. \label{eq:capital_dynamics}
            \end{equation}

            В отличие от капитала, переход по потреблению задаётся неявно. Подставляя доходность капитала $R_{t+1} = R(K_{t+1}, C_{t+1})$ в уравнение Эйлера, получаем соотношение
            \begin{equation}
            C_{t+1} = \beta \bigl(1 + R(K_{t+1}, C_{t+1}) - \delta \bigr) C_t, \label{eq:euler_implicit}
            \end{equation}
            где $R(K_{t+1}, C_{t+1})$ зависит от $K_{t+1}$ и $C_{t+1}$ через эндогенный труд $L_{t+1}$ и выпуск $Y_{t+1}$.

            Таким образом, исходная модель сводится к двумерной динамической системе относительно переменных состояния $K_t$ и управления $C_t$. Переход по капиталу (\ref{eq:capital_dynamics}) вычисляется явно, тогда как для определения $C_{t+1}$ требуется на каждом шаге решать нелинейное уравнение (\ref{eq:euler_implicit}) относительно неизвестного $C_{t+1}$.
        \section{Итоговая динамическая система}

            Численная реализация модели основана на последовательном вычислении переменных каждого периода в строго определённом порядке. Пусть в начале периода $t$ известны капитал $K_t$ и потребление $C_t$. Тогда выполняются следующие действия.
            
            \begin{equation}
            \begin{cases}
            L_t = \left( \frac{(1-\alpha)A}{\phi} \cdot \frac{K_t^{\alpha}}{C_t} \right)^{\frac{1}{\psi+\alpha}}, \\[4pt]
            Y_t = A K_t^{\alpha} L_t^{1-\alpha}, \\[4pt]
            R_t = \alpha \frac{Y_t}{K_t}, \\[4pt]
            w_t = (1-\alpha) \frac{Y_t}{L_t}  \\[4pt]
            I_t = Y_t - C_t \\[4pt]
            K_{t+1} = Y_t - C_t + (1-\delta) K_t, \\[4pt]
            C_{t+1} = \beta \bigl(1 + R(K_{t+1}, C_{t+1}) - \delta \bigr) C_t.
            \end{cases}
            \end{equation}

            Система замыкается начальными условиями: $K_0$ задано, $C_0$ подбирается так, чтобы траектория сходилась к стационарному состоянию. Таким образом, мы имеем двумерную динамическую систему с явным переходом по капиталу и неявным — по потреблению.
    % --- Глава 2 ---

    \chapter{Обзор численных методов решения задач оптимального управления в макроэкономике}

        \section{Подходы к получению необходимых условий оптимальности}

            В задачах динамического общего равновесия, к которым относится модель реального делового цикла, экономические агенты принимают решения на бесконечном временном горизонте, максимизируя межвременную функцию полезности при технологических и ресурсных ограничениях. Для вывода условий, описывающих оптимальную траекторию, используются три основных подхода.

            Первый подход, наиболее распространённый для дискретных моделей, основан на методе Лагранжа с динамическими множителями (множителями Куна--Таккера). Составляется функция Лагранжа, включающая целевую функцию и все ограничения, после чего условия первого порядка приводят к системе разностных уравнений, включая уравнение Эйлера для потребления и условие трансверсальности.

            Второй подход, широко применяемый в моделях непрерывного времени, опирается на принцип максимума Понтрягина. Вводится функция Гамильтона, и оптимальное управление находится из условия её максимизации, а сопряжённые переменные удовлетворяют дифференциальным уравнениям.

            Третий подход — рекурсивный, основанный на принципе оптимальности Беллмана. Вместо последовательности управлений ищется функция ценности (value function), удовлетворяющая функциональному уравнению Беллмана. Этот подход особенно полезен при численном решении методом итераций по функции ценности.

            Независимо от способа вывода, получаемая система уравнений (условия оптимальности) в общем случае является нелинейной и требует применения численных методов. В макроэкономической литературе сложилось разделение на локальные и глобальные методы решения таких систем \cite{judd, heermaussner}.

        \section{Локальные методы}

            Локальные методы основаны на аппроксимации исходной нелинейной модели в окрестности стационарного состояния. Наиболее распространёнными являются линеаризация и лог-линеаризация, когда все переменные заменяются их логарифмическими отклонениями от равновесных значений. В результате исходная система сводится к линейной системе разностных уравнений, которая может быть решена аналитически или методом спектрального разложения.

            Для дискретных моделей рациональных ожиданий, к которым относится RBC-модель, ключевую роль играет метод Бланшара--Кана \cite{blanchardkahn}. Этот метод позволяет определить, существует ли единственная устойчивая траектория, исходя из расположения собственных значений матрицы, описывающей линеаризованную динамику. Если число собственных значений, лежащих внутри единичного круга, совпадает с числом предопределённых переменных (переменных состояния), то решение существует и единственно.

            Преимущества локальных методов заключаются в их относительной простоте, низкой вычислительной стоимости и возможности аналитического исследования устойчивости. Однако они имеют существенный недостаток: точность аппроксимации быстро убывает при удалении от стационарного состояния. В случае сильных шоков или значительных начальных отклонений линеаризация может давать качественно неверные результаты \cite{heermaussner}. Кроме того, локальные методы не позволяют исследовать глобальные свойства модели, такие как множественность равновесий или эффекты, связанные с ограничениями на переменные.
        
        \section{Глобальные методы}

            Глобальные методы не предполагают, что траектория находится в малой окрестности стационарного состояния. Они строят решение на всём пространстве состояний, хотя, как правило, ценой существенно большего объёма вычислений. Среди глобальных методов можно выделить следующие.

            \textbf{Метод стрельбы (shooting)} является одним из наиболее естественных для детерминированных моделей с седловой динамикой. Идея метода состоит в сведении краевой задачи (известны начальный капитал и условие трансверсальности на бесконечности) к задаче подбора недостающего начального значения потребления. Траектория прошивается (интегрируется) вперёд, и начальное потребление корректируется до тех пор, пока на конечном горизонте траектория не попадёт в достаточно малую окрестность стационарного состояния. Различают простую стрельбу (single shooting), когда интегрируется вся траектория целиком, и множественную стрельбу (multiple shooting), когда горизонт разбивается на участки для улучшения устойчивости счёта.

            \textbf{Сеточные методы} (например, метод итерации по функции ценности) дискретизируют пространство состояний и на каждом узле сетки вычисляют функцию ценности, решая уравнение Беллмана. Эти методы гарантируют нахождение глобального решения, но страдают от «проклятия размерности»: объём вычислений экспоненциально растёт с числом переменных состояния \cite{judd}. Для двух-трёх переменных состояния они ещё применимы, но для более сложных моделей требуют специальных приёмов (например, аппроксимации функции ценности сплайнами или нейронными сетями).

            \textbf{Проекционные методы} приближают искомые функции (политику потребления или функцию ценности) в виде разложения по конечному базису (например, полиномам Чебышёва или сплайнам). Коэффициенты разложения находятся из условия, чтобы аппроксимация удовлетворяла уравнению Беллмана или условиям оптимальности в выбранных узлах (метод коллокации). Проекционные методы часто оказываются более эффективными, чем сеточные, при гладких решениях \cite{judd, miranda}.
        \section{Выбор метода для настоящей работы}

            В данной работе исследуется детерминированная RBC-модель с эндогенным предложением труда и неявным уравнением Эйлера для потребления. Численное решение должно обеспечить построение переходной траектории при заданном начальном капитале и гарантировать сходимость к стационарному состоянию.

            
            Метод стрельбы является естественным выбором по следующим причинам:
            
            \begin{itemize}
                \item модель детерминирована и имеет седловую динамику, что делает метод стрельбы хорошо обоснованным;
                \item размерность пространства состояний невелика (одна переменная состояния $K_t$), что позволяет использовать простую стрельбу без потери устойчивости;
                \item метод непосредственно работает с исходными нелинейными уравнениями, не требуя линеаризации;
                \item относительно прост в реализации и даёт интерпретируемые результаты.
            \end{itemize}

            Наряду с методом стрельбы в работе проводится анализ стационарного состояния и свойств линеаризованной системы, что позволяет сопоставить глобальное численное решение с предсказаниями локальной теории. Это даёт дополнительную проверку корректности результатов и позволяет оценить степень нелинейности модели при выбранной калибровке.

            В качестве альтернативы можно было бы применить метод итерации по функции ценности на сетке, однако при наличии труда как эндогенной переменной это потребовало бы значительных вычислительных ресурсов без явного преимущества в точности для детерминированного случая. Поэтому основным выбранным методом является метод стрельбы, который подробно описан в следующей главе.
    % --- Глава 3 ---

    \chapter{Исследование и построение решения задачи}
        \section{Аналитическое исследование стационарного состояния}

            Стационарным состоянием (steady state) модели называется такая траектория, на которой все переменные остаются постоянными во времени. Иными словами, выполняются условия:
            \[
            K_{t+1}=K_t=K^*, \qquad 
            C_{t+1}=C_t=C^*, \qquad 
            L_{t+1}=L_t=L^*,
            \]
            и, как следствие, $Y_{t+1}=Y_t=Y^*$, $R_{t+1}=R_t=R^*$, $w_{t+1}=w_t=w^*$, $I_{t+1}=I_t=I^*$.

            Вывод аналитических выражений для стационарных значений осуществляется последовательно, начиная с уравнения Эйлера.


            Уравнение Эйлера для потребления в стационарном режиме принимает вид:
            \begin{equation}
            1 = \beta \bigl(1 + R^* - \delta\bigr).
            \end{equation}
            Отсюда немедленно получаем:
            \begin{equation}
            R^* = \frac{1}{\beta} - 1 + \delta. \tag{3.1}
            \end{equation}
            Таким образом, стационарная доходность капитала полностью определяется параметрами дисконтирования $\beta$ и выбытия $\delta$.


            Закон движения капитала в стационарном состоянии даёт:
            \[
            K^* = Y^* - C^* + (1-\delta)K^* \quad\Longrightarrow\quad Y^* = C^* + \delta K^*.
            \]

            С другой стороны, из определения доходности капитала $R^* = \alpha Y^*/K^*$ имеем $Y^* = \frac{R^*}{\alpha}K^*$. Приравнивая два выражения для $Y^*$, находим:
            \[
            \frac{R^*}{\alpha}K^* = C^* + \delta K^* \quad\Longrightarrow\quad 
            \frac{R^*}{\alpha} = \frac{C^*}{K^*} + \delta.
            \]

            Введём обозначение $v = C^*/K^*$ для отношения стационарного потребления к капиталу. Тогда:
            \begin{equation}
            v = \frac{C^*}{K^*} = \frac{R^*}{\alpha} - \delta. \tag{3.2}
            \end{equation}


            Для нахождения $L^*$ используем условие внутрипериодной оптимальности, связывающее заработную плату, потребление и труд:
            \[
            \phi C^* (L^*)^{\psi} = w^*,
            \]
            а также выражение для заработной платы через предельный продукт труда:
            \[
            w^* = (1-\alpha)\frac{Y^*}{L^*}.
            \]

            Подставляя $Y^* = \frac{R^*}{\alpha}K^*$ и $C^* = v K^*$, получаем:
            \[
            \phi v K^* (L^*)^{\psi} = (1-\alpha)\frac{R^*}{\alpha}\frac{K^*}{L^*}.
            \]

            Сокращая $K^*$ (положительный множитель) и умножая обе части на $L^*$, приходим к уравнению:
            \[
            \phi v (L^*)^{\psi+1} = \frac{1-\alpha}{\alpha}R^*.
            \]

            Отсюда выражается стационарный уровень труда:
            \begin{equation}
            L^* = \left( \frac{(1-\alpha)R^*}{\alpha \phi v} \right)^{\frac{1}{\psi+1}}. \tag{3.3}
            \end{equation}


            Теперь найдём $K^*$, используя производственную функцию Кобба–Дугласа:
            \[
            Y^* = A (K^*)^{\alpha} (L^*)^{1-\alpha}.
            \]

            С другой стороны, из $R^* = \alpha Y^*/K^*$ имеем $Y^* = \frac{R^*}{\alpha}K^*$. Приравнивая, получаем:
            \[
            \frac{R^*}{\alpha}K^* = A (K^*)^{\alpha} (L^*)^{1-\alpha}.
            \]

            Разрешая относительно $K^*$:
            \[
            (K^*)^{1-\alpha} = \frac{\alpha A}{R^*} (L^*)^{1-\alpha}
            \quad\Longrightarrow\quad
            K^* = \left( \frac{\alpha A}{R^*} \right)^{\frac{1}{1-\alpha}} L^*.
            \]

            Более удобная форма (явно выделяющая $L^*$):
            \begin{equation}
            K^* = \left( \frac{\alpha A (L^*)^{1-\alpha}}{R^*} \right)^{\frac{1}{1-\alpha}}. \tag{3.4}
            \end{equation}

            Наконец, стационарное потребление находится по определению $v$:
            \begin{equation}
            C^* = v K^*. \tag{3.5}
            \end{equation}


            Таким образом, стационарное состояние модели полностью описывается следующей системой аналитических выражений, вычисляемых в строгой последовательности:

            \[
            \begin{aligned}
            R^* &= \frac{1}{\beta} - 1 + \delta, \\[4pt]
            v &= \frac{R^*}{\alpha} - \delta, \\[4pt]
            L^* &= \left( \frac{(1-\alpha)R^*}{\alpha \phi v} \right)^{\frac{1}{\psi+1}}, \\[4pt]
            K^* &= \left( \frac{\alpha A (L^*)^{1-\alpha}}{R^*} \right)^{\frac{1}{1-\alpha}}, \\[4pt]
            C^* &= v K^*, \\[4pt]
            Y^* &= A (K^*)^{\alpha} (L^*)^{1-\alpha} = \frac{R^*}{\alpha} K^*, \\[4pt]
            I^* &= \delta K^*, \\[4pt]
            w^* &= (1-\alpha)\frac{Y^*}{L^*}.
            \end{aligned}
            \]

            Все величины выражаются исключительно через экзогенные параметры модели $\beta$, $\alpha$, $\delta$, $\phi$, $\psi$, $A$, что позволяет проводить численные расчёты и анализировать чувствительность равновесия к изменению параметров.
            
        \section{Численные значения стационарного состояния}

            Для проведения вычислительных экспериментов необходимо задать численные значения экзогенных параметров модели. В работе используется калибровка, широко применяемая в литературе по моделям реального делового цикла \cite{kingrebello, cooley}. Выбор параметров отражает следующие экономические соображения:

            \begin{itemize}
                \item Коэффициент дисконтирования $\beta = 0.95$ соответствует годовой реальной процентной ставке около $5\%$, что типично для развитых экономик.
                \item Доля капитала в доходе $\alpha = 0.33$ является стандартной оценкой для экономики США и большинства стран ОЭСР.
                \item Параметр совокупной факторной производительности $A = 0.95$ выбран таким образом, чтобы стационарный выпуск был сопоставим с масштабом капитала.
                \item Норма выбытия капитала $\delta = 0.10$ означает, что ежегодно выбывает $10\%$ основного капитала, что соответствует сроку службы оборудования около 10 лет.
                \item Параметр дизполезности труда $\phi = 0.75$ определяет интенсивность неприятия труда домохозяйством.
                \item Эластичность предложения труда по Фришу $\psi = 0.35$ лежит в диапазоне эмпирических оценок (обычно от 0.2 до 0.5).
                \item Параметр $\sigma = 1.0$ соответствует логарифмической функции полезности по потреблению (межвременная эластичность замещения равна единице). Хотя в выписанной системе уравнений этот параметр явно не входит, он используется в исходной постановке задачи полезности домохозяйства.
            \end{itemize}

            Таким образом, набор параметров имеет вид:
            \[
            \beta = 0.95,\quad \alpha = 0.33,\quad A = 0.95,\quad \sigma = 1.0,\quad \delta = 0.1,\quad \phi = 0.75,\quad \psi = 0.35.
            \]


            Подставляя указанные параметры в аналитические формулы (3.1)–(3.5), полученные в предыдущем разделе, находим численные значения стационарных переменных. Вычисления проводились с двойной точностью; ниже приведены результаты с округлением до десяти значащих цифр.

            \begin{align}
            R^* &= 0.15263157894736837, \\
            K^* &= 3.2258188530080525, \\
            C^* &= 1.1694236447986124, \\
            L^* &= 1.101749946608162, \\
            Y^* &= 1.4920055300994175, \\
            I^* &= 0.3225818853008053.
            \end{align}


            Полученное стационарное значение труда $L^* \approx 1.12$ превышает единицу. Это означает, что в стационарном режиме домохозяйство тратит на работу более $100\%$ своего времени, что экономически некорректно, поскольку труд не может превышать доступного временного endowment (обычно нормированного на единицу). Формально данное значение не препятствует численному решению системы, однако указывает на необходимость корректировки калибровки.

            Причина превышения единицы кроется в выборе параметров $\phi$, $\psi$ и $A$. Снижение параметра дезполезности труда $\phi$ (меньшая тягость труда) или увеличение совокупной факторной производительности $A$ ведёт к росту $L^*$. Для получения стационарного труда в интервале $(0, 1)$ следует либо увеличить $\phi$, либо уменьшить $A$, либо выбрать иные значения $\alpha$ и $\beta$. В рамках данной работы аномальное значение $L^*$ не является критическим, поскольку метод стрельбы остаётся работоспособным, а качественное поведение переходных траекторий сохраняется. Тем не менее, в перспективе требуется более тщательная калибровка, возможно, с использованием метода моментов или байесовской оценки.

            Для контроля правильности вычислений можно проверить выполнение основных соотношений стационарного состояния:
            \[
            Y^* - C^* - \delta K^* \approx 1.49200 - 1.16942 - 0.1*3.22581 \approx 0.00000 \ (\text{в пределах точности}),
            \]
            \[
            R^* - \alpha\frac{Y^*}{K^*} = 0.15263 - 0.33 \times \frac{1.49200}{3.22581} \approx 0.15263 - 0.15263 = 0.
            \]
            Таким образом, вычисленные значения согласованы с уравнениями модели.
        
    \chapter{Метод стрельбы}
        
        Для нахождения устойчивой переходной траектории используется метод стрельбы. Его идея состоит в следующем: при фиксированном начальном значении капитала $K_0$ подбирается такое начальное значение потребления $C_0$, при котором траектория системы на конечном горизонте максимально приближается к стационарному состоянию.

        На каждом шаге алгоритма выполняются следующие действия:
        \begin{enumerate}[label=\arabic*)]
            \item по заданным $K_t$ и $C_t$ вычисляется $L_t$;
            \item определяется выпуск $Y_t$ и инвестиции $I_t=Y_t-C_t$;
            \item по уравнению накопления вычисляется капитал $K_{t+1}$;
            \item численно решается нелинейное уравнение относительно $C_{t+1}$:
            \begin{equation}
            C_{t+1}=\beta\bigl(1+R(K_{t+1},C_{t+1})-\delta\bigr)C_t;
            \end{equation}
            \item процесс повторяется до достижения заданного горизонта моделирования.
        \end{enumerate}

        Подбор начального потребления осуществляется по минимуму функционала, измеряющего отклонение терминального состояния от стационарной точки.
            
        \section{Построение траектории при фиксированном \texorpdfstring{$C_0$}{C0}}

            Пусть заданы начальные значения $K_0$ и $C_0$. Для каждого периода $t=0,1,\dots,T-1$ траектория строится итерационно. Сначала по внутрипериодному условию вычисляется труд:
            \[
            L_t=\left(\frac{(1-\alpha)A}{\phi}\frac{K_t^\alpha}{C_t}\right)^{\frac{1}{\psi+\alpha}}.
            \]
            Затем определяются выпуск и связанные с ним величины:
            \(Y_t=A K_t^\alpha L_t^{1-\alpha}\),
            \(R_t=\alpha\frac{Y_t}{K_t}\),
            \(I_t=Y_t-C_t\).
            После этого капитал следующего периода находится из закона накопления:
            \(K_{t+1}=Y_t-C_t+(1-\delta)K_t\).

            Значение $C_{t+1}$ определяется из уравнения Эйлера и находится как корень одномерного нелинейного уравнения
            \[
            f(c)=c-\beta\bigl(1+R(K_{t+1},c)-\delta\bigr)C_t=0,
            \]
            где \(c=C_{t+1}\).
            При этом функция доходности на шаге \(t+1\) задается через
            \(Y_{t+1}=A K_{t+1}^\alpha L_{t+1}^{1-\alpha}\),
            \(L_{t+1}=\left(\frac{(1-\alpha)A}{\phi}\frac{K_{t+1}^\alpha}{c}\right)^{\frac{1}{\psi+\alpha}}\),
            \(R_{t+1}=\alpha\frac{Y_{t+1}}{K_{t+1}}\).
            После нахождения \(C_{t+1}\) выполняется переход к периоду \(t+1\), и процедура повторяется до горизонта \(T\).
        \section{Метод Ньютона для нахождения \texorpdfstring{$C_{t+1}$}{C(t+1)}}

            После вычисления \(K_{t+1}\) значение \(C_{t+1}\) определяется из неявного уравнения Эйлера:
            \[
            f(c)=c-\beta\bigl(1+R(K_{t+1},c)-\delta\bigr)C_t=0,\qquad c=C_{t+1}.
            \]
            Для фиксированного \(K=K_{t+1}\) используем обозначения
            \(L(K,c)=\left(\frac{(1-\alpha)A}{\phi}\frac{K^\alpha}{c}\right)^{\frac{1}{\psi+\alpha}}\),
            \(Y(K,c)=A K^\alpha L(K,c)^{1-\alpha}\),
            \(R(K,c)=\alpha\frac{Y(K,c)}{K}\).

            Производная \(f'(c)\) вычисляется аналитически. Из
            \(\frac{\partial L}{\partial c}=-\frac{1}{\psi+\alpha}\frac{L}{c}\),
            \(\frac{\partial Y}{\partial c}=-\frac{1-\alpha}{\psi+\alpha}\frac{Y}{c}\)
            получаем
            \(\frac{\partial R}{\partial c}=-\frac{1-\alpha}{\psi+\alpha}\frac{R}{c}\), а значит
            \[
            f'(c)=1-\beta C_t\frac{\partial R(K,c)}{\partial c}
            =1+\beta C_t\frac{1-\alpha}{\psi+\alpha}\frac{R(K,c)}{c}.
            \]

            Тогда итерации Ньютона записываются в виде
            \[
            c^{(n+1)}=
            c^{(n)}-\frac{f(c^{(n)})}{f'(c^{(n)})}
            =
            c^{(n)}-
            \frac{c^{(n)}-\beta\bigl(1+R(K_{t+1},c^{(n)})-\delta\bigr)C_t}
            {1+\beta C_t \frac{1-\alpha}{\psi+\alpha}\frac{R(K_{t+1},c^{(n)})}{c^{(n)}}}.
            \]
            После достижения критерия сходимости принимается \(C_{t+1}=c^{(n)}\).
        
        \section{Псевдокод метода стрельбы}
            \begin{enumerate}[label=\arabic*)]
                \item Задать параметры модели $(\alpha,\beta,\delta,\phi,\psi,A)$ и вычислить стационарное состояние $(K^*,C^*,L^*)$.
                \item Зафиксировать $K_0$ и начальный интервал для потребления: $C_0\in[C_{\min},C_{\max}]$.
                \item Повторять внешнюю бисекцию по $C_0$, пока $|C_{\max}-C_{\min}|\ge\varepsilon_{\text{outer}}$.
                \item На каждой итерации положить $C_0=\frac{C_{\min}+C_{\max}}{2}$ и построить траекторию на горизонте $t=0,1,\dots,T-1$.
                \item Для каждого $t$ вычислить
                \(L_t=\left(\frac{(1-\alpha)A}{\phi}\frac{K_t^\alpha}{C_t}\right)^{\frac{1}{\psi+\alpha}}\),
                \(Y_t=A K_t^\alpha L_t^{1-\alpha}\),
                \(K_{t+1}=Y_t-C_t+(1-\delta)K_t\),
                затем найти $C_{t+1}$ методом Ньютона из
                \(f(c)=c-\beta\bigl(1+R(K_{t+1},c)-\delta\bigr)C_t=0\).
                \item По завершении траектории вычислить терминальное отклонение $g(C_0)=K_T-K^*$.
                \item Обновить интервал: если $g(C_0)>0$, то $C_{\min}=C_0$, иначе $C_{\max}=C_0$.
                \item После сходимости внешней бисекции принять найденное $C_0^*$ и построить итоговую saddle-path траекторию $(K_t,C_t,L_t)$.
            \end{enumerate}



    % --- Глава 4 ---

    \chapter{Описание практической части решения детерминированной модели}
        Численная реализация выполнена на языке Python с использованием библиотек \texttt{NumPy}, \texttt{SciPy}, \texttt{Matplotlib} и \texttt{Plotly}. Библиотека \texttt{SciPy} использовалась для решения одномерного нелинейного уравнения и подбора стартового значения $C_0$. Библиотека \texttt{Matplotlib} применялась для построения статических графиков, а \texttt{Plotly} --- для интерактивной визуализации фазовых траекторий.

        Алгоритм реализован в виде набора функций, каждая из которых отвечает за вычисление отдельных компонент модели: труда, выпуска, доходности капитала, следующего значения капитала, а также решения неявного уравнения для потребления. Поверх этих функций построена процедура генерации траектории на конечном горизонте и процедура оптимального подбора начального значения потребления.

        \section{Эксперимент 1: старт ниже стационарного уровня капитала}

            В первом эксперименте рассматривалось начальное условие
            \begin{equation}
            K_0=0.3K^*.
            \end{equation}

            Методом стрельбы было найдено значение
            \begin{equation}
            C_0\approx 0.702956.
            \end{equation}

            При таком выборе начального потребления траектория переменных $K_t$, $C_t$, $L_t$ ,$Y_t$ сходится к стационарному состоянию. Экономически это соответствует ситуации, когда экономика стартует с недостаточным объемом капитала и постепенно выходит на долгосрочное равновесие.

                На рис.~\ref{fig:K_low_1},~\ref{fig:K_low_2},~\ref{fig:K_low_3d} показана динамика основных показателей при старте ниже стационарного уровня ($K_0=0.3K^*.$).
                \begin{figure}[H]
                    \centering
                    \includegraphics[width=\textwidth]{images/rbc_trajectory_2d_K0_1.0991_T_45.png}
                    \caption{Переходная динамика $K_t,Y_t,C_t,L_t$ при $K_0<K^*$}
                    \label{fig:K_low_1}
                \end{figure}

                \begin{figure}[H]
                    \centering
                    \includegraphics[width=\textwidth]{images/rbc_trajectory_2d_additional_K0_1.0991_T_45.png}
                    \caption{Переходная динамика $R_t,w_t,I_t$ при $K_0<K^*$}
                    \label{fig:K_low_2}
                \end{figure}


                \begin{figure}[H]
                    \centering
                    \includegraphics[width=\textwidth]{images/rbc_trajectory_3d_1.0991.png}
                    \caption{Трёхмерная фазовая траектория модели при $K_0<K^*$}
                    \label{fig:K_low_3d}
                \end{figure}
        \section{Эксперимент 2: старт выше стационарного уровня капитала}

            Во втором эксперименте использовалось начальное условие
            \begin{equation}
            K_0=10K^*.
            \end{equation}

            Подбором методом стрельбы было найдено значение
            \begin{equation}
            C_0\approx 4.542463.
            \end{equation}

            В этом случае экономика стартует с избыточным объемом капитала по сравнению со стационарным уровнем. Численные расчеты показывают, что и в данной ситуации траектория постепенно приближается к steady state.


                На рис.~\ref{fig:K_high_1},~\ref{fig:K_high_2},~\ref{fig:K_high_3d} показана динамика основных показателей при старте выше стационарного уровня ($K_0=10K^*.$).
                \begin{figure}[H]
                    \centering
                    \includegraphics[width=\textwidth]{images/rbc_trajectory_2d_K0_36.6358_T_56.png}
                    \caption{Переходная динамика $K_t,Y_t,C_t,L_t$ при $K_0>K^*$}
                    \label{fig:K_high_1}
                \end{figure}

                \begin{figure}[H]
                    \centering
                    \includegraphics[width=\textwidth]{images/rbc_trajectory_2d_additional_K0_36.6358_T_56.png}
                    \caption{Переходная динамика $R_t,w_t,I_t$ при $K_0>K^*$}
                    \label{fig:K_high_2}
                \end{figure}


                \begin{figure}[H]
                    \centering
                    \includegraphics[width=\textwidth]{images/rbc_trajectory_3d_36.6358.png}
                    \caption{Трёхмерная фазовая траектория модели при $K_0>K^*$}
                    \label{fig:K_high_3d}
                \end{figure}

        \section{Оценка вычислительной схемы}

            Проведенные вычислительные эксперименты показывают, что метод стрельбы позволяет устойчиво строить переходные траектории для рассматриваемой модели. Точность результата зависит от способа решения неявного уравнения для $C_{t+1}$, длины горизонта моделирования и выбора критерия близости к стационарному состоянию. При этом учет труда делает модель более содержательной, но одновременно усложняет нелинейную структуру динамики.

    \chapter{Исследование модели со стохастическими шоками}

        В детерминированной модели, рассмотренной в предыдущих главах, все переменные эволюционируют по заданным законам при фиксированных параметрах и начальных условиях. Однако реальные экономические данные демонстрируют колебания, которые невозможно объяснить только детерминированной динамикой. Для воспроизведения таких флуктуаций в модель вводятся случайные возмущения — \textbf{стохастические шоки}. Эти шоки интерпретируются как неожиданные изменения в технологии (шоки совокупной факторной производительности), денежно-кредитной политике, предпочтениях потребителей или государственных расходах.

        В настоящей работе рассматриваются два типа шоков:
        \begin{itemize}
            \item \textbf{Шок совокупной факторной производительности (TFP)} — отражает непредвиденные изменения в технологии производства, эффективности использования капитала и труда. Положительный шок TFP увеличивает выпуск при тех же затратах факторов.
            \item \textbf{Денежно-кредитный (монетарный) шок} — моделирует неожиданное изменение в правиле денежно-кредитной политики (например, отклонение процентной ставки от целевого значения).
        \end{itemize}

        Включение шоков превращает детерминированную модель в модель динамического стохастического общего равновесия (DSGE). Решением такой модели является не одна траектория, а распределение возможных траекторий, определяемое вероятностными свойствами шоков.

        \section{Математическая постановка стохастической модели}

        В стохастической версии модели некоторые параметры или экзогенные переменные становятся случайными процессами. В данной работе используется следующая спецификация:

        \begin{enumerate}
            \item Производственная функция включает стохастическую совокупную факторную производительность $A_t$ (в коде обозначена как \texttt{tfp}):
            \[
            Y_t = A_t K_t^{\alpha} L_t^{1-\alpha}.
            \]
            
            \item Логарифм TFP следует авторегрессионному процессу первого порядка AR(1):
            \[
            \ln A_t = \rho_{A} \ln A_{t-1} + \varepsilon_t^{A}, \quad \varepsilon_t^{A} \sim \mathcal{N}(0, \sigma_A^2),
            \]
            где $\rho_A \in (0,1)$ — коэффициент автокорреляции, $\varepsilon_t^{A}$ — белый шум (инновация) с дисперсией $\sigma_A^2$.
            
            \item Денежно-кредитное правило (правило Тейлора) в логарифмической форме:
            \[
            \ln R_t^b = \gamma \ln R_{t-1}^b + (1-\gamma)\left[ \ln R^{ss} + \tau \left( \ln \pi_{t+1} - \ln \pi^{ss} \right) \right] + \varepsilon_t^{mon}, 
            \]
            где $R_t^b$ — номинальная процентная ставка, $\pi_t$ — инфляция, $\gamma$ — степень сглаживания ставки, $\tau$ — чувствительность к ожидаемой инфляции, $\varepsilon_t^{mon}$ — монетарный шок.
        \end{enumerate}

        Остальные уравнения модели (условия оптимальности домохозяйств и фирм, бюджетные ограничения, тождества) остаются такими же, как в детерминированном случае, но теперь все переменные становятся случайными величинами, а ожидания будущих значений играют ключевую роль.

        \section{Методы решения стохастических DSGE-моделей}

        Точное аналитическое решение нелинейной стохастической системы с рациональными ожиданиями, как правило, недоступно. Поэтому применяются численные методы, которые можно разделить на три класса.

        \paragraph*{Локальная линеаризация (первого порядка)}
        Наиболее распространённый подход для DSGE-моделей — линеаризация условий первого порядка в окрестности детерминированного стационарного состояния. Исходная нелинейная система заменяется линейной системой разностных уравнений относительно отклонений переменных от стационарных значений (обычно в логарифмах). Преимущества: относительная простота, возможность применения методов решения линейных моделей рациональных ожиданий (например, метод Бланшара–Кана \cite{blanchardkahn}), быстрота вычислений. Недостатки: точность только в малой окрестности стационарного состояния; не учитываются нелинейные эффекты (например, премия за риск).

        \paragraph*{Методы высших порядков}
        Линеаризация второго или третьего порядка позволяет уловить некоторые нелинейные свойства (например, асимметрию реакции на положительные и отрицательные шоки, зависимость волатильности от состояния экономики). Однако вычислительная сложность резко возрастает.

        \paragraph*{Глобальные методы}
        К ним относятся проекционные методы, методы итерации по функции ценности на сетке и методы, основанные на аппроксимации полиномами высоких степеней. Эти методы не требуют близости к стационарному состоянию, но их применение к моделям с несколькими переменными состояния сопряжено с «проклятием размерности».

        В данной работе используется \textbf{линеаризация первого порядка}, реализованная в программном пакете Dynare, поскольку она даёт достаточную точность для анализа импульсных переходных функций и моментов при умеренной волатильности шоков.
        
        \section{Алгоритм решения линеаризованной модели}

        После линеаризации система приводится к виду:
        \[
        \Gamma_0 \mathbf{x}_t = \Gamma_1 \mathbf{x}_{t-1} + \Psi \boldsymbol{\varepsilon}_t + \Pi \eta_t,
        \]
        где $\mathbf{x}_t$ — вектор всех переменных модели (в отклонениях), $\boldsymbol{\varepsilon}_t$ — вектор структурных шоков, $\eta_t$ — ошибки ожиданий (рациональные ожидания). Метод Бланшара–Кана позволяет найти решение в форме:
        \[
        \mathbf{x}_t = \Phi \mathbf{x}_{t-1} + \Phi_{\varepsilon} \boldsymbol{\varepsilon}_t.
        \]

        Dynare автоматически выполняет:
        \begin{enumerate}
            \item символьное вычисление стационарного состояния;
            \item линеаризацию (или лог-линеаризацию) уравнений модели;
            \item решение методом Бланшара–Кана;
            \item вычисление импульсных переходных функций (IRF) — реакции переменных на единичный шок;
            \item расчёт теоретических моментов (стандартных отклонений, корреляций) и декомпозицию дисперсии по шокам.
        \end{enumerate}

        \section{Интерпретация импульсных переходных функций}

        Импульсная переходная функция показывает, как переменные модели реагируют во времени на однократное шоковое воздействие (обычно величиной в одно стандартное отклонение). В работе анализируются отклики ключевых макроэкономических переменных ($Y_t$, $C_t$, $I_t$, $L_t$, $\pi_t$, $R_t$) на шоки TFP и монетарный шок. Положительный шок TFP, как правило, приводит к росту выпуска, потребления и инвестиций, а также к временному снижению инфляции (если цены гибкие). Монетарный шок (неожиданное повышение процентной ставки) вызывает спад выпуска и инфляции в соответствии с каналом трансмиссии денежно-кредитной политики.

        \section{Краткое описание реализации в Dynare}

        Программная реализация стохастической модели выполнена в среде Octave с использованием пакета Dynare (версия 6.5). Исходный файл \texttt{RBC\_RPMP.mod} содержит:
        \begin{itemize}
            \item объявление переменных и параметров;
            \item описание модели в виде системы нелинейных уравнений с рациональными ожиданиями;
            \item задание стационарного состояния в блоке \texttt{steady\_state\_model};
            \item спецификацию стохастических шоков (AR(1)-процессы для TFP и монетарного правила);
            \item команду \texttt{stoch\_simul} для линеаризации и расчёта импульсных откликов.
        \end{itemize}

        Результаты моделирования (графики IRF, моменты, декомпозиция дисперсии) получены в виде EPS-файлов и числовых таблиц. В следующем разделе представлен анализ этих результатов.

    \chapter{Результаты стохастического моделирования: импульсные переходные функции}

        В данном разделе представлены результаты численного решения стохастической модели, полученные с помощью пакета Dynare. Для анализа реакции экономики на внешние возмущения были рассчитаны импульсные переходные функции (IRF) на два типа шоков: шок совокупной факторной производительности (TFP) и монетарный шок. Все графики сохранены в формате png.

        \section{Импульсные отклики на шок TFP}

            На рис.~\ref{fig:Y_tfp}–\ref{fig:L_tfp} показаны реакции основных переменных на положительный шок совокупной факторной производительности.

            \begin{figure}[H]
                \centering
                \includegraphics[width=0.8\textwidth]{graphs_png_improved/Y_tfp_shock.png}
                \caption{Импульсный отклик выпуска $Y$ на шок TFP.}
                \label{fig:Y_tfp}
            \end{figure}

            \begin{figure}[H]
                \centering
                \includegraphics[width=0.8\textwidth]{graphs_png_improved/C_tfp_shock.png}
                \caption{Импульсный отклик потребления $C$ на шок TFP.}
                \label{fig:C_tfp}
            \end{figure}

            \begin{figure}[H]
                \centering
                \includegraphics[width=0.8\textwidth]{graphs_png_improved/K_tfp_shock.png}
                \caption{Импульсный отклик капитала $K$ на шок TFP.}
                \label{fig:K_tfp}
            \end{figure}

            \begin{figure}[H]
                \centering
                \includegraphics[width=0.8\textwidth]{graphs_png_improved/L_tfp_shock.png}
                \caption{Импульсный отклик труда $L$ на шок TFP.}
                \label{fig:L_tfp}
            \end{figure}
        

            Положительный шок TFP приводит к росту выпуска, потребления и инвестиций, что согласуется с теоретическими ожиданиями. При этом инвестиции реагируют наиболее сильно, поскольку они сглаживают потребление во времени. Труд (занятость) может незначительно снижаться в краткосрочном периоде из-за эффекта дохода (домохозяйства предпочитают больше отдыхать при росте производительности). Инфляция временно снижается, а центральный банк реагирует на это снижением номинальной процентной ставки, что дополнительно стимулирует экономику.
        
        \section{Импульсные отклики на монетарный шок}

        
            \begin{figure}[H]
                \centering
                \includegraphics[width=0.8\textwidth]{graphs_png_improved/Y_mon_shock.png}
                \caption{Импульсный отклик выпуска $Y$ на монетарный шок.}
                \label{fig:Y_mon}
            \end{figure}

            \begin{figure}[H]
                \centering
                \includegraphics[width=0.8\textwidth]{graphs_png_improved/C_mon_shock.png}
                \caption{Импульсный отклик потребления $C$ на монетарный шок.}
                \label{fig:C_mon}
            \end{figure}

            \begin{figure}[H]
                \centering
                \includegraphics[width=0.8\textwidth]{graphs_png_improved/K_mon_shock.png}
                \caption{Импульсный отклик капитала $K$ на монетарный шок.}
                \label{fig:K_mon}
            \end{figure}

            \begin{figure}[H]
                \centering
                \includegraphics[width=0.8\textwidth]{graphs_png_improved/L_mon_shock.png}
                \caption{Импульсный отклик труда $L$ на монетарный шок.}
                \label{fig:L_mon}
            \end{figure}

            Рост номинальной процентной ставки ведёт к снижению инвестиций и потребления (из-за удорожания кредита), что вызывает спад выпуска. Инфляция также снижается, так как падает совокупный спрос. Труд сокращается вслед за выпуском. Реальная процентная ставка (номинальная ставка минус ожидаемая инфляция) в краткосрочном периоде повышается, что соответствует стандартному каналу трансмиссии денежно-кредитной политики.

    % --- Выводы (по требованию) ---

    \chapter*{Выводы}
        \addcontentsline{toc}{chapter}{Выводы}
        В результате выполнения работы получены следующие основные результаты:
        \begin{enumerate}[label=\arabic*)]
            \item аналитически выведены формулы для стационарного состояния детерминированной RBC-модели с трудом;
            \item проведён обзор локальных и глобальных численных методов решения задач оптимального управления в макроэкономике;
            \item реализован численный алгоритм метода стрельбы с учётом неявного уравнения Эйлера для потребления;
            \item выполнены вычислительные эксперименты для двух начальных уровней капитала (ниже и выше стационарного);
            \item построены графики переходных траекторий и фазовые портреты, подтверждающие сходимость к стационарному состоянию.
        \end{enumerate}

    % --- Заключение (по желанию) ---

    \chapter*{Заключение}
        \addcontentsline{toc}{chapter}{Заключение}
        Проведённое исследование показало, что численные результаты чувствительны к параметрам, определяющим предложение труда. Перспективным направлением дальнейшей работы является более точная калибровка параметров на основе экономически интерпретируемых ограничений, а также расширение модели за счёт включения технологических шоков, применения локальных методов решения и анализа стохастической версии RBC-модели.

    \clearpage
    \phantomsection

    \addcontentsline{toc}{chapter}{Список литературы}
        \begin{thebibliography}{99}

        \bibitem{blanchardkahn}
        Blanchard O., Kahn C. The Solution of Linear Difference Models under Rational Expectations. // Econometrica, 1980, Vol. 48, No. 5, pp. 1305--1311.
        
        \bibitem{romer}
        Бессонов В. А. Проблемы построения производственных функций в российской переходной экономике. — Институт экономики переходного периода, 2002.

        \bibitem{acemoglu}
        Acemoglu D. Introduction to Modern Economic Growth. — Princeton: Princeton University Press, 2009.

        \bibitem{heermaussner}
        Heer B., Maussner A. Dynamic General Equilibrium Modeling: Computational Methods and Applications. — Berlin: Springer, 2009.

        \bibitem{ljungqvist}
        Ljungqvist L., Sargent T. Recursive Macroeconomic Theory. — Cambridge, MA: MIT Press, 2018.

        \bibitem{miranda}
        Miranda M., Fackler P. Applied Computational Economics and Finance. — Cambridge, MA: MIT Press, 2002.

        \bibitem{judd}
        Judd K. Numerical Methods in Economics. — Cambridge, MA: MIT Press, 1998.

        \bibitem{kingrebello}
        King R., Rebelo S. Resuscitating Real Business Cycles. // Handbook of Macroeconomics. — Elsevier, 1999.

        \bibitem{cooley}
        Cooley T. (ed.) Frontiers of Business Cycle Research. — Princeton University Press, 1995.

        \end{thebibliography}

\end{document}